\section{Introduction}\label{sec:intro}

Meibomian gland dysfunction (MGD) is among the most common causes of dry eye disease\cite{foulks2003,villani2020latest}. Infrared/near-infrared meibography visualises gland dropout, tortuosity, and shortening non-invasively, and grading systems such as Meiboscore map dropout to clinical severity\cite{wise2012meibography,ngo2014repeat}. Manual grading is observer-dependent and only moderately repeatable. Pixel-wise segmentation of glands and eyelid margin is therefore a prerequisite for quantitative area, length, tortuosity, and R\_loss.

This task differs from generic semantic segmentation in three respects. First, targets are extremely thin: at $640\times1280$, a gland is 3--8 pixels wide. Any stride-16 backbone reduces it below one pixel; fragmentation appears before Dice drops, and is exposed mainly by HD95 and skeleton continuity. Second, contrast is low and cross-device variation is large. Third, data are scarce: about one thousand paired cases cannot train a large Transformer from scratch, making pretraining and strong regularisation necessary\cite{saha2022automated,setu2021deep}.

Public data have changed the field. MGD-1K provides 1000 LipiView II infrared images, gland and eyelid masks, demographics for 320 subjects, and six expert Meiboscore rounds per image\cite{saha2022automated}; recent multicentre work and MGD-1K baselines show that public data can support reproducible comparison\cite{li2025multicenter,yu2025quantitative}. Public data, however, do not automatically yield comparability. We identify three remaining problems.

\textbf{(i) Patient-level leakage.} The official \officialTestImages-case Meiboscore test list contains \officialTestPatients{} patients, while \officialLeakImages{} of the remaining 800 images come from the same patients. Image-level random splits let the same patient's upper/lower eyelids and left/right eyes cross training and test. Even with a patient-level split, using the official 200 directly still leaves \officialLeakPatients/\officialTestPatients{} test patients in the training set. Leakage is not ``using more data'': it lets subject identity become a memorised feature and systematically inflates generalisation.

\textbf{(ii) Hidden effects of tiling.} Tile height, overlap, threshold, and TTA are often buried in implementation detail or omitted. With identical weights and only tile height changed, we observe Dice moving from 0.370 to 0.655, while full-image inference drops to 0.593. The same model can therefore be narrated as strong or weak, and readers cannot tell whether the difference comes from architecture or inference configuration.

\textbf{(iii) The gap between pixel metrics and clinical endpoints.} Dice, HD95, and clDice describe pixels and topology, not Meiboscore or gland dropout. Six expert grading rounds let us directly compare the segmentation$\to$R\_loss$\to$Meiboscore chain with clinical grades, and to test upper/lower eyelid differences. Prior work often excludes the lower eyelid because of incomplete exposure, which hides real clinical difficulty.

This paper therefore treats comparability and leakage audit as the primary contribution, not a novel module. Our contributions are:

\begin{enumerate}[label=(\arabic*)]
\item \textbf{A patient-level public benchmark.} On \nTotal{} deduplicated MGD-1K pairs we build a patient-independent split: \nTrain/\nVal/\nTest{} images (\nTrainPatients/\nValPatients/\nTestPatients{} patients), and release the split file, the file-list signature, and the assignment signature, so that ``same file list, different assignment'' is no longer untraceable.
\item \textbf{Leakage audit and quantification.} We report the patient composition of the official \officialTestImages-case test list (\officialTestPatients{} patients), the \officialLeakImages/800{} image-level leakage, and the \officialLeakPatients/\officialTestPatients{} patient-level leakage; we construct a \cleanOfficialImages-image/\cleanOfficialPatients-patient clean-official subset and compare patient-level versus leaky protocols.
\item \textbf{A five-route comparison under one protocol.} UNet, Attention U-Net, WRCANet, MGD-PriorNet, and MGD-Former are trained and evaluated with the same preprocessing, the same patient-level split, the same loss, and the same tiling and threshold. WRCANet did not converge within the 8\,GB budget, so the main table ranks the four comparable routes and reports WRCANet separately as a memory/time feasibility counterexample. The four comparable routes are baselines and ablation branches, not failed attempts being replaced.
\item \textbf{Tiling as a first-class variable and Multi-Tile Ensemble (MTE).} We report tile sensitivity from 128 to 512 and full image with identical weights, and propose MTE, which averages probabilities from 256/384/512 tiles; we report its HD95 boundary improvement (Dice did not improve) and compare it with 8$\times$ TTA in accuracy and inference cost.
\item \textbf{Uncertainty-guided review.} We compare multi-tile variance, predictive entropy, and 8$\times$ TTA by AUROC for failure cases (Dice $<0.7$) and report their limitations; variance is only weak-to-moderate and TTA8 suffers from label--signal coupling, so neither is a standalone triage rule.
\item \textbf{Clinical grading and subgroup analysis.} We report accuracy, MAE, and quadratic weighted kappa between segmentation-derived R\_loss/Meiboscore and six-round expert grades, and report Dice/clDice separately for upper and lower eyelids.
\end{enumerate}

No accuracy claim is fabricated. All Dice/IoU/HD95/clDice, AUROC, and clinical agreement values will be injected by script after unified training and independent testing; tiling sensitivity and VRAM/throughput are engineering measurements labelled with their conditions. Section~\ref{sec:method} describes the method and training recipe, Section~\ref{sec:engineering} turns patient-level splitting, leakage audit, and MTE into methodological objects, Section~\ref{sec:experiments} gives the experiment matrix, Section~\ref{sec:results} presents the result structure and discussion, and Section~\ref{sec:limitations} lists limitations.
